\chapter{Ringraziamenti}

Questo rappresenta per me il capitolo più difficile da scrivere: rendere a parole e trasmettere un concetto su un qualcosa su cui studi ogni giorno, sbattendoci la testa fino ad vedere quel barlume di illuminazione che ti fa capire effettivamente cosa sta succedendo, è molto più facile, almeno a mio avviso, piuttosto che mettere su carta quello che vorresti dire a tutte quelle persone che, silenziosamente o meno, ti hanno accompagnato in questo lungo percorso. \\
Vorrei ringraziare, innanzitutto, il mio relatore, il professore Leonardo Gualtieri, per la sua disponibilità ad essere il mio relatore, oltre ad aver rappresentato, almeno per me, un docente fantastico, di cui ho avuto l'immensa fortuna di averlo sia come esercitatore del corso di Fisica 1 e sia come docente titolare del corso di Meccanica Quantistica, che ritengo essere il corso che ho apprezzato di più, sia per argomenti trattati e sia per la trattazione, nel percorso triennale di Fisica. \\
Subito dopo, vorrei ringraziare la mia famiglia, che mi ha supportato costantemente in questi anni, senza mai dubitare delle mie capacità: ringrazio i miei genitori, Antonietta (perdonami mamma se ti chiamo col nome intero) ed Enrico, per tutto il supporto morale che mi avete dato in questi anni; in particolar modo la mamma, che fra le mie numerose sfide e le sue numerose sfide e responsabilità, è sempre stata lì, nei momenti di bisogno, ad ascoltarmi nei (numerosi) momenti di sconforto e disperazione senza mai farmi pesare nulla. Successivamente, non posso non ringraziare la zia Adele e lo zio Maurizio, che mi hanno sempre supportato quasi come se fossi figlio loro, e mia cugina Giada, che ha sempre creduto nel suo cuginetto (anche se, purtroppo, i ``compiti'' di matematica che faccio adesso non puoi più correggerli facilmente come facevi un tempo). Infine, ringrazio la nonna Enrica, a cui ho realizzato il grande sogno di far vedere entrambi i nipoti laureati, e ringrazio sempre per la grande mano che ha dato a me e mia mamma in questi tre anni. \\
Non sarei riuscito a farcela senza le grandi amicizie che ho fatto in questi anni, di cui, ovviamente, devo ringraziare i ``quarratini'' OG, ovvero Carmine, Dario e Ale, che ormai conosco da tempo immemore, che hanno avuto modo di sperimentare la mia grandissima pazzia già prima di Fisica e hanno visto il suo andamento esponenziale in questi anni di ``cattività pisana'' in ogni chiamata su Discord. Ringrazio (fra virgolette anche se spero non me ne abbiate a male) i quarratini ``aggiunti'', Matteo, Sam, Michael e Tommaso, qua presenti, che ho conosciuto dopo, ma non di meno (in particolar modo Matteo che a Ferragosto di ritorno da Lavagna mi ha salvato la vita) che sto piano piano corrompendo verso la lega dei campioni \\
Ringrazio i miei coinquilini Anna e Matteo, con cui ho iniziato questo percorso $3$ anni fa, anche se il nostro ``cammino'' insieme è iniziato ben prima, quando sono arrivato il primo giorno di terza superiore al liceo scientifico. Adesso siamo tutti separati ahimé, ma, nonostante le normali discussioni che ci possono essere fra coinquilini, mi avete sempre sostenuto nei numerosi momenti del bisogno, spesso dinanzi ad un bel pasto dato da una ``minoranza'' (nuovo nome ufficiale di Just Eat).  \\
E ora iniziano i ringraziamenti dei ''pisani``: ringrazio Ismaele Bartoluccio, il mio daddy topologico, che è certamente una fra le persone più ``ganze'' che io abbia mai conosciuto (alla fine sono convenuto per questo termine, l'unico che può forse rendere su carta la grandezza dello spirito libero e indipendente di Isma): in questi anni mi ha preso sotto la sua ala, spingendomi nello studio (disperatissimo) delle aree della matematica più affascinanti ed esotiche per un fisico; oltre, nell'ultimo anno, ad avermi dato un supporto morale per tutta una serie di vicende della mia vita che, in questa sede, è meglio non raccontare, ma che Ismaele sa bene. Ringrazio poi Riccardo, che mi ha supportato e sopportato ogni giorno, ogni ora e ogni minuto di questa triennale, subendo i miei grandissimi danni psicologici ``alla 6pek'' dovuti alla mia pazzia (forse saranno stati questi che ti hanno convinto a fare Teorica insieme, piuttosto che Astrofisica ahah). Ringrazio Alessandro Terranova, che rimarrà il mio solo e unico presidentessimo, per avermi sostenuto più e più volte nei momenti in cui ero più vulnerabile, oltre a farmi conoscere il meraviglioso mondo di AI-SF (non Associazione Italiana Studio del Fegato obv) che, ad entrambi, rappresenta una grande famiglia estesa (sia per le grandissimi emozioni che ci procura e sia per i grandi momenti di tribolazione che ci ha procurato). Ringrazio il mitico Francesco Zeno Costanzo (o più brevemente il Gen) per avermi preso, a suon di insulti per le interlinee assenti ``alla Bolo'', sotto la sua ala pythonica e per avermi fatto capire quanto sia importante scrivere una stringa di spiegazione sia fondamentale (anche per la più banale delle funzioni) e per intarmi la \emph{toplane} (caro si scherza). Ringrazio Alberto Otera (il mio goat) per avermi fatto anche lui da psicologo dopo le lezioni pazze ed incredibili di Gruppi (mettersi accanto forse, come dicevi te, non era forse una buona idea) e di MQA, oltre a graziarmi con le sue mosse divine su League of Legends. Ringrazio Francesco Angelo Fabiano Antonacci per aver fatto parte in questo grandissimo esperimento sociale che consisteva nel porre, in un gruppo di laboratorio, la persona più teorica del dipartimento con, forse, la persona più sperimentale del dipartimento e per non aver mai perso la sua grandissima capacità nel dare i nomi nei file utili a laboratorio (ricordo a tutti l'esplicativo \emph{four\_int\_der\_train2.py}) e ringrazio Benedetta Santoro  
